% Pursuing Stacks (cotes 134-2 to 134-8): archive cote 134-8, site page 26 (PDF page 27), typescript page 580.
% typescript. Draft compared: ps7.tex:1201-1263.
% AI-transcribed by Claude Opus 5.5 (claude-opus-5-5, adaptive thinking, effort high), Message Batches API, batch ps4 (msgbatch_012Aoc1NuskATPsyBeKSYK7J), 2026-11-12; not yet checked by eye.
\scanpage{134-8/26}
\mspage{580}
common middle, hence is contractible. Thus the $R^{0}\hand{|}f\hand{|}_{*}$ of a constant
sheaf on the disc $\hand{|}I\hand{|}$ is by no means constant on the open\hand{,} dimension one
stratum of $\hand{|}J\hand{|}$ .

These examples bring\xout{s} to my mind that for any map $f: I \to J$ between
finite ordered sets, \xout{satisfying the} possibly submitted to the mild
restriction that the sets (13) should \over{\unsure{b}}{b}e non-empty, the homotopy types
of those \xout{set} ordered sets\inserted{(13)} (for the order relation induced by I) should
be exactly the homotopy types of the fibers of the maps \xout{fxxxxxxxxxx}\inserted{(12) between strata.}
Thus asphericity of these ordered sets should express no more, no less
than the con\hand{t}ractibility of those fibers.\deleted{(} This latter condition is
exactly what is needed in order to insure stability by $R\hand{|}f\hand{|}_{*}$ of the
notion of complexes of sheaves compatible with the stratifications.

28.10.

138. Yesterday and the day before\hand{,} I got involved \xout{again} in that
unforeseen digression\deleted{\unsure{,}} around the foreboding of a ``six operations
duality formalism'' in (Cat), and suitable notions for smoothness and
(more important still) of properness for a map in (Cat). This digres-%
sion wholly convinced me that the usual duality formalism should hold
in (Cat) too, Working this out in full detail should be a most pleasant
task indeed, and presumably the best, or even the only way for gaining
complete mastery of the cohomology formalism within (Cat) - or, what
more or less amounts to the same, for topoi which admit sufficiently
many projective objects. In the previous two sections, I referred to
such duality formalism as\inserted{one} concerned with sheaves of k-modules, for any
fixed ring k - but from the example of étale duality \xout{say, it} for sche-%
mes say, it is likely that instead of fixing a ring k, we may as
well take objects A in (Cat) endowed with an arbitrary sheaf of rings
$\underline{O}_{A}$ (which we'll only have to suppose commutative, \over{\unsure{wh}}{wh}en concerned with
the two\hand{\unsure{,}}internal operations $\overset{\hand{L}}{\otimes}$ and $R\underline{Hom}$ ), and taking maps of such
ringed objects as the basic maps. In the present context, the\inserted{usual} ``six
operations'' in duality theory will be enriched however by still another
\margin{! $\times$} one, namely $Lf_{\hand{!}}$ , \xout{valid}\inserted{defined} for any map f between such ringed objects\deleted{\unsure{,}}
\hand{(}not to be confused with the $Rf_{\hand{!}}$ operation\hand{,} \deleted{(}defined only under sui-%
table finiteness asum\hand{\illegible{1}}ptions, such as ``properness'', for the underly-%
ing map in (Cat)), whose relationship to the other operations should
be understood and \deleted{included}\inserted{added} \over{in}{to} the standard duality formulaire. One

\margin{26}
